\chapter{Kernel Limits for Wide Networks: NNGP and NTK}\label{ch:kernels}

The first regime considered in this chapter is the infinite-width initialization limit, realized here at an unlimited hidden width $N \in \hNinf$. Neal showed that a one-hidden-layer network induces a Gaussian process at random initialization~\cite{neal1996priors}; Lee et al.\ and Matthews et al.\ extended this NNGP description to finite-depth networks~\cite{lee2018deep,matthews2018gaussian}.

We then turn to a logically distinct training-time regime. In the NTK, or lazy-training, regime at unlimited width $N \in \hNinf$, the empirical tangent kernel at each standard training time concentrates infinitesimally around a deterministic standard kernel $K$ that does not depend on time, so the standard-part output dynamics reduce to kernel gradient flow~\cite{jacot2018neural,lee2019wide}.

\section{Neural Network Gaussian Processes (NNGP)}\label{sec:ch2-nngp}

We first consider a one-hidden-layer network of unlimited hidden width $N \in \hNinf$ over a standard input dimension $n_0 \in \N$. Write $[N] := \{1, \dots, N\} \subset \hN$, and for each hidden neuron $i \in [N]$ set $h_i(x) := n_0^{-1/2}\ip{w_i}{x}$. At hyperfinite width $N$, the network's scalar output on a standard input $x \in \R^{n_0}$ is the internal hyperfinite sum
\begin{equation}\label{eq:ch2-1layer-net}
f(x; \theta) = \frac{1}{\sqrt{N}} \sum_{i=1}^{N} a_i \, \ns{\varphi}(h_i(x)), \qquad x \in \R^{n_0},
\end{equation}
where the internal parameters are $\theta = \{(a_i, w_i)\}_{i=1}^{N}$ defined on an internal probability space $(\Omega, \mathcal{A}, \hP)$ equipped with its countably additive Loeb probability measure $\Loeb$~\cite{loeb1975conversion,cutland2000loeb}. At initialization we assume internal Gaussian weights
\begin{equation}\label{eq:ch2-init-weights}
a_i \stackrel{\text{i.i.d.}}{\sim} \hNcal(0, 1), \quad
w_i \stackrel{\text{i.i.d.}}{\sim} \hNcal(0, I_{n_0}), \quad
\text{independent across $i$ and between $a_i$ and $w_i$.}
\end{equation}

For fixed standard inputs $x^1, \dots, x^m \in \R^{n_0}$ (with $m \in \N$ standard), let $\mathcal{F} := \ns{\sigma}(h_i(x^\alpha) : i \in [N],\, \alpha \in [m]) \subset \mathcal{A}$ be the internal $\sigma$-algebra generated by the hidden preactivations. Thus $\mathcal{F}$ is generated by the hidden neurons at the inputs under consideration and does not include the output weights. Since the internal Gaussian output weights are internally independent of $\mathcal{F}$, by the Transfer Principle~\cite{robinson1966nonstandard} the output vector is exactly internally conditionally Gaussian:
\begin{equation}\label{eq:ch2-cond-gauss}
\bigl(f(x^\alpha; \theta)\bigr)_{\alpha=1}^{m} \,\Big|\, \mathcal{F} \sim \hNcal\bigl(0, \Phi^{(N)}\bigr), \qquad
\Phi^{(N),\alpha\beta} := \frac{1}{N} \sum_{i=1}^{N} \ns{\varphi}\bigl(h_i(x^\alpha)\bigr) \, \ns{\varphi}\bigl(h_i(x^\beta)\bigr).
\end{equation}
Thus no conditional central limit theorem is needed in the Gaussian-output-weight setting. This output-layer identity uses only internal Gaussian readout weights independent of the hidden neurons and will be reused below. If the relevant standard second moments are finite, the hyperfinite (Loeb) strong law of large numbers gives, $\Loeb$-almost surely,
\begin{equation}\label{eq:ch2-slln-phi}
\Phi^{(N),\alpha\beta} \approx \Phi(x^\alpha, x^\beta), \qquad
\Phi(x, x') := \E\left[ \varphi\!\left(\frac{\ip{w}{x}}{\sqrt{n_0}}\right) \varphi\!\left(\frac{\ip{w}{x'}}{\sqrt{n_0}}\right) \right], \quad w \sim \mathcal{N}(0, I_{n_0}).
\end{equation}
Taking internal conditional characteristic functions in~\eqref{eq:ch2-cond-gauss} and passing to standard parts via Loeb bounded convergence therefore yields the exact standard-part Gaussian law
\begin{equation}\label{eq:ch2-f-shadow-law}
\st\bigl(f(x^1), \dots, f(x^m)\bigr) \sim \mathcal{N}\Bigl(0, \bigl[\Phi(x^\alpha, x^\beta)\bigr]_{\alpha,\beta=1}^{m}\Bigr) \qquad \text{under } \Loeb.
\end{equation}
For centered, unit-variance non-Gaussian output weights, the analogous conclusion instead follows from a hyperfinite multivariate central limit theorem.

\begin{exercise}[Gaussianity of Linear Combinations]\label{ex:ch2-linear-comb}
Show that for any standard $m \in \N$ and any standard coefficients $u \in \R^m$, $\st\bigl(\sum_{\alpha=1}^{m} u_\alpha f(x^\alpha)\bigr) \sim \mathcal{N}(0, u^\top \Phi u)$ under $\Loeb$, with $\Phi = [\Phi(x^\alpha, x^\beta)]_{\alpha,\beta=1}^{m}$.
\end{exercise}

The calculation above establishes the finite-dimensional standard-part distributions, which identify the shadow Gaussian process on every fixed standard collection of inputs. Realizing the shadow of the random function itself on continuous function spaces additionally requires S-continuity of the internal sample paths.

\begin{definition}[Gaussian Process]\label{def:ch2-gp}
A random function $f : \R^{n_0} \to \R$ is a Gaussian process with mean $\mu$ and covariance kernel $\Phi$ (notation: $f \sim \GP(\mu, \Phi)$) if
\begin{enumerate}[label=(\roman*)]
\item $\Phi$ is symmetric positive semidefinite, i.e.\ for any $r \in \N$ and any $x^1, \dots, x^r \in \R^{n_0}$, the matrix $[\Phi(x^\alpha, x^\beta)]_{\alpha,\beta=1}^{r}$ is PSD.
\item For any $r \in \N$ and any $x^1, \dots, x^r \in \R^{n_0}$,
\begin{equation}\label{eq:ch2-gp-fdd}
\bigl(f(x^1), \dots, f(x^r)\bigr) \sim \mathcal{N}\Bigl(\bigl(\mu(x^1), \dots, \mu(x^r)\bigr), \bigl[\Phi(x^\alpha, x^\beta)\bigr]_{\alpha,\beta=1}^{r}\Bigr).
\end{equation}
\end{enumerate}
\end{definition}

\begin{theorem}[Finite-Dimensional NNGP Limit at Initialization~\cite{neal1996priors}]\label{thm:ch2-nngp-fdd}
Assume $\varphi : \R \to \R$ is Borel measurable and, for every $x \in \R^{n_0}$,
\begin{equation}\label{eq:ch2-phi-l2}
\E\left[ \varphi\!\left(\frac{\ip{w}{x}}{\sqrt{n_0}}\right)^2 \right] < \infty, \qquad w \sim \mathcal{N}(0, I_{n_0}).
\end{equation}
Then, at any unlimited hidden width $N \in \hNinf$ and for every fixed standard finite collection $x^1, \dots, x^m \in \R^{n_0}$, the vector $\bigl(f(x^\alpha; \theta)\bigr)_{\alpha=1}^{m}$ belongs to $\Lim^m$ $\Loeb$-almost surely and its standard part satisfies
\begin{equation}\label{eq:ch2-nngp-thm-law}
\st\Bigl(\bigl(f(x^\alpha; \theta)\bigr)_{\alpha=1}^{m}\Bigr) \sim \mathcal{N}\Bigl(0, \bigl[\Phi(x^\alpha, x^\beta)\bigr]_{\alpha,\beta=1}^{m}\Bigr) \qquad \text{under } \Loeb,
\end{equation}
where $\Phi$ is the kernel in~\eqref{eq:ch2-slln-phi}. Equivalently, the finite-dimensional standard-part distributions coincide with those of a centered Gaussian process with covariance kernel $\Phi$.
\end{theorem}

\begin{proof}
By~\eqref{eq:ch2-phi-l2} and Transfer, the internally i.i.d.\ products $\ns{\varphi}(h_i(x^\alpha))\ns{\varphi}(h_i(x^\beta))$ have standard-finite internal expectation $\Phi(x^\alpha, x^\beta)$, so the Loeb strong law of large numbers in~\eqref{eq:ch2-slln-phi} gives $\Phi^{(N),\alpha\beta} \approx \Phi(x^\alpha, x^\beta)$ $\Loeb$-almost surely. By~\eqref{eq:ch2-cond-gauss}, $\hE[\norm{(f(x^\alpha))_{\alpha=1}^{m}}_2^2 \mid \mathcal{F}] = \Tr(\Phi^{(N)}) \in \Lim$ $\Loeb$-almost surely, so $(f(x^\alpha; \theta))_{\alpha=1}^{m} \in \Lim^m$ $\Loeb$-almost surely. For any standard $u \in \R^m$, taking internal expectations of $\hE[\ns{\exp}(i\sum_{\alpha=1}^{m} u_\alpha f(x^\alpha; \theta)) \mid \mathcal{F}] = \ns{\exp}(-\frac{1}{2} u^\top \Phi^{(N)} u)$ and applying Loeb bounded convergence gives $\EL[\exp(i\sum_{\alpha=1}^{m} u_\alpha \st(f(x^\alpha; \theta)))] = \exp(-\frac{1}{2} u^\top \Phi u)$, proving~\eqref{eq:ch2-nngp-thm-law}. Moreover, for every $u \in \R^m$,
\begin{equation}\label{eq:ch2-phi-psd}
\sum_{\alpha,\beta=1}^{m} u_\alpha u_\beta \Phi(x^\alpha, x^\beta)
= \E\left[ \left( \sum_{\alpha=1}^{m} u_\alpha \varphi\!\left(\frac{\ip{w}{x^\alpha}}{\sqrt{n_0}}\right) \right)^2 \right] \ge 0.
\end{equation}
Thus $\Phi$ is positive semidefinite and defines the stated centered Gaussian process.
\end{proof}

\begin{remark}[Function-Space Strengthening]\label{rem:ch2-func-space}
The theorem characterizes finite-dimensional standard-part distributions, which is the notion needed for a fixed finite dataset. Lifting $f$ to a continuous standard Gaussian process in $C(K)$ on a standard compact set $K \subset \R^{n_0}$ additionally requires $\Loeb$-almost-sure S-continuity of the internal sample paths on $\ns{K}$ (that is, $f(x) \in \Lim$ for all $x \in \ns{K}$ and $f(x) \approx f(x')$ whenever $x, x' \in \ns{K}$ satisfy $x \approx x'$). For example, a globally H\"older activation yields the uniform monadic H\"older increment bound $\hE[\abs{f(x) - f(x')}^p] \le C_{p,K}\norm{x - x'}^{\gamma p}$ proved in Appendix~A.2; the nonstandard Kolmogorov--Chentsov criterion then guarantees S-continuity of $f$ on $\ns{K}$ $\Loeb$-almost surely, so $f^\circ(\st(x)) := \st(f(x))$ defines a continuous standard Gaussian process $f^\circ \sim \GP(0, \Phi|_{K \times K})$ on $K$ satisfying $\sup_{x \in \ns{K}} \abs{f(x) - \ns{f^\circ}(x)} \approx 0$ $\Loeb$-almost surely.
\end{remark}

For some activations the kernel $\Phi$ admits a closed form. For ReLU, $\varphi(u) = \max\{u, 0\}$, the relevant Gaussian expectations have explicit expressions (the Cho--Saul / arc-cosine kernel).

\begin{proposition}[Cho--Saul Kernel for ReLU]\label{prop:ch2-cho-saul}
Let $(h^\alpha, h^\beta)$ be a centered Gaussian pair with
\begin{equation}\label{eq:ch2-cho-saul-cov}
\E\bigl[(h^\alpha)^2\bigr] = \Phi^{\alpha\alpha}, \qquad
\E\bigl[(h^\beta)^2\bigr] = \Phi^{\beta\beta}, \qquad
\E\bigl[h^\alpha h^\beta\bigr] = \Phi^{\alpha\beta}.
\end{equation}
Assume $\Phi^{\alpha\alpha}, \Phi^{\beta\beta} > 0$ and define the correlation
\begin{equation}\label{eq:ch2-cho-saul-rho}
\rho := \frac{\Phi^{\alpha\beta}}{\sqrt{\Phi^{\alpha\alpha}\Phi^{\beta\beta}}} \in [-1, 1].
\end{equation}
Then
\begin{align}
\E\bigl[\varphi(h^\alpha)\varphi(h^\beta)\bigr] &= \frac{\sqrt{\Phi^{\alpha\alpha}\Phi^{\beta\beta}}}{2\pi}\left( \sqrt{1 - \rho^2} + \rho\left(\frac{\pi}{2} + \arcsin(\rho)\right) \right), \label{eq:ch2-cho-saul-fwd} \\
\E\bigl[\varphi'(h^\alpha)\varphi'(h^\beta)\bigr] &= \Prob(h^\alpha > 0,\, h^\beta > 0) = \frac{1}{4} + \frac{1}{2\pi}\arcsin(\rho). \label{eq:ch2-cho-saul-deriv}
\end{align}
\end{proposition}

\begin{proof}
By positive homogeneity of ReLU, it is enough to work with a jointly standard Gaussian pair $(z_1, z_2)$ with correlation $\rho$ and restore the marginal variances at the end. First suppose $\abs{\rho} < 1$. Let $w_1, w_2$ be independent standard Gaussians and realize $z_1 = w_1$ and $z_2 = \rho w_1 + \sqrt{1 - \rho^2}\, w_2$. Let $f$ and $F$ denote the standard Gaussian density and distribution function, and write $p(\rho) := \Prob(z_1 > 0,\, z_2 > 0)$. Conditioning on $w_1$ gives
\begin{equation}\label{eq:ch2-orthant-int}
p(\rho) = \int_0^\infty f(u) F\!\left(\frac{\rho u}{\sqrt{1 - \rho^2}}\right) du.
\end{equation}
Differentiation under the integral is justified on compact subintervals of $(-1, 1)$ by Gaussian domination, and yields
\begin{equation}\label{eq:ch2-orthant-deriv}
p'(\rho) = \frac{1}{2\pi(1 - \rho^2)^{3/2}} \int_0^\infty u \exp\!\left(-\frac{u^2}{2(1 - \rho^2)}\right) du = \frac{1}{2\pi\sqrt{1 - \rho^2}}.
\end{equation}
Since $p(0) = 1/4$, integration gives the orthant probability stated in~\eqref{eq:ch2-cho-saul-deriv}. Because $\varphi'(z) = \1_{\{z > 0\}}$ almost everywhere, this proves the derivative-kernel formula for $\abs{\rho} < 1$.

We now compute the expectation in~\eqref{eq:ch2-cho-saul-fwd}. Let $f_\rho$ be the joint density of $(z_1, z_2)$. For dummy coordinates $u, v$, direct differentiation gives $u f_\rho(u, v) = -\partial_u f_\rho(u, v) - \rho\, \partial_v f_\rho(u, v)$. Multiply this identity by $v$ and integrate over the positive quadrant. Gaussian decay removes the boundary terms at infinity, and the factor $v$ removes the boundary term at $v = 0$. The $u$-integration leaves the boundary at $u = 0$; integration by parts in $v$ leaves $p(\rho)$. Consequently,
\begin{equation}\label{eq:ch2-relu-ibp}
\E\bigl[\varphi(z_1)\varphi(z_2)\bigr] = \int_0^\infty v f_\rho(0, v)\, dv + \rho p(\rho) = \frac{\sqrt{1 - \rho^2}}{2\pi} + \rho p(\rho).
\end{equation}
The boundary integral is an elementary one-dimensional Gaussian integral. Substituting the orthant probability already obtained gives the bracketed expression in~\eqref{eq:ch2-cho-saul-fwd} for the standard pair.

Finally, $(h^\alpha, h^\beta) \eqd (\sqrt{\Phi^{\alpha\alpha}}\, z_1, \sqrt{\Phi^{\beta\beta}}\, z_2)$. Positive homogeneity therefore multiplies the ReLU expectation by $\sqrt{\Phi^{\alpha\alpha}\Phi^{\beta\beta}}$, while positive rescaling leaves the signs unchanged. This proves both stated formulas for $\abs{\rho} < 1$. The cases $\rho = \pm 1$ follow from the same coupling: bounded convergence applies to the indicator product, and the Lipschitz property of ReLU together with Gaussian second moments gives continuity of the ReLU-product expectation.
\end{proof}

\begin{remark}\label{rem:ch2-cho-saul}
Formulas~\eqref{eq:ch2-cho-saul-fwd}--\eqref{eq:ch2-cho-saul-deriv} go back to Cho and Saul's arc-cosine kernels~\cite{cho2009kernel} (and by the Transfer Principle hold verbatim internally for centered $\hNcal(0, \Sigma)$ pairs with $\Sigma^{\alpha\alpha}, \Sigma^{\beta\beta} \in \hR_{>0}$). They are also the standard closed forms used to compute the limiting NNGP/NTK kernels for ReLU networks.
\end{remark}

\begin{remark}[Bayesian View]\label{rem:ch2-bayesian}
At unlimited width $N \in \hNinf$, random initialization induces a standard-part GP prior $f^\circ \sim \GP(0, \Phi)$ under $\Loeb$. This motivates GP-style Bayesian inference as an approximation to Bayesian neural nets.
\end{remark}

\section{Extension to Deep Networks: Kernel Recursion}\label{sec:ch2-deep-nngp}

To extend the NNGP construction beyond one hidden layer, fix a standard depth $d \in \N$ with $d \ge 2$ and an unlimited width $N \in \hNinf$ for each hidden layer. For standard data points $\{x^\alpha\}_{\alpha=1}^{m} \subset \R^{n_0}$ (with $m, n_0 \in \N$ standard), define an internal multilayer perceptron (MLP) by
\begin{align}
h_1^\alpha &:= \frac{1}{\sqrt{n_0}} W_0 x^\alpha, & W_0 &\in \hR^{N \times n_0}, \label{eq:ch2-mlp-h1} \\
h_{\ell+1}^\alpha &:= \frac{1}{\sqrt{N}} W_\ell \, \ns{\varphi}(h_\ell^\alpha), & W_\ell &\in \hR^{N \times N}, \quad \ell = 1, \dots, d-1, \label{eq:ch2-mlp-hell} \\
f(x^\alpha; \theta) &:= \frac{1}{\sqrt{N}} W_d \, \ns{\varphi}(h_d^\alpha), & W_d &\in \hR^{1 \times N}. \label{eq:ch2-mlp-out}
\end{align}
All weight entries are internally i.i.d.\ $\hNcal(0, 1)$ on $(\Omega, \mathcal{A}, \hP)$ and internally independent across layers, with $\ns{\varphi}$ applied coordinatewise.

We will repeatedly use the following standard Gaussian facts, all of which carry over verbatim to arbitrary hyperfinite dimensions $N, M \in \hN$ (with $\hNcal$ in place of $\mathcal{N}$) by the Transfer Principle.

\begin{proposition}[Linear Images of Gaussians]\label{prop:ch2-gauss-linear}
If $g \sim \mathcal{N}(\mu, \Sigma)$ in $\R^n$ and $A \in \R^{m \times n}$, then $Ag \sim \mathcal{N}(A\mu, A\Sigma A^\top)$.
\end{proposition}

\begin{proposition}[Inner Products with a Standard Gaussian]\label{prop:ch2-gauss-ip}
If $g \sim \mathcal{N}(0, I_n)$ and $u, v \in \R^n$, then
\begin{equation}\label{eq:ch2-gauss-ip}
\begin{pmatrix} \ip{g}{u} \\ \ip{g}{v} \end{pmatrix}
\sim \mathcal{N}\!\left(0, \begin{pmatrix} \norm{u}^2 & \ip{u}{v} \\ \ip{u}{v} & \norm{v}^2 \end{pmatrix}\right).
\end{equation}
\end{proposition}

\begin{proposition}[Multiplying by a Gaussian Matrix]\label{prop:ch2-gauss-matvec}
Let $W \in \R^{n \times n}$ have i.i.d.\ entries $W_{ij} \sim \mathcal{N}(0, 1)$ and let $u, v \in \R^n$ be deterministic. Then
\begin{equation}\label{eq:ch2-gauss-matvec}
\begin{pmatrix} Wu \\ Wv \end{pmatrix}
\sim \mathcal{N}\!\left(0, \begin{pmatrix} \norm{u}^2 I_n & \ip{u}{v} I_n \\ \ip{u}{v} I_n & \norm{v}^2 I_n \end{pmatrix}\right)
= \mathcal{N}\!\left(0, \begin{pmatrix} \norm{u}^2 & \ip{u}{v} \\ \ip{u}{v} & \norm{v}^2 \end{pmatrix} \otimes I_n\right).
\end{equation}
\end{proposition}

\begin{definition}[Kronecker Product]\label{def:ch2-kronecker}
If $A = [a_{ij}] \in \R^{p \times q}$ and $B \in \R^{r \times s}$ (or internally for $p, q, r, s \in \hN$), the Kronecker product $A \otimes B \in \R^{pr \times qs}$ is defined blockwise by
\begin{equation}\label{eq:ch2-kronecker}
A \otimes B := [a_{ij} B]_{i,j}.
\end{equation}
\end{definition}

Define the deterministic standard input Gram matrix
\begin{equation}\label{eq:ch2-phi0}
\Phi_0^{\alpha\beta} := \frac{1}{n_0} \ip{x^\alpha}{x^\beta},
\end{equation}
and set $\Phi_0^{(N)} := \Phi_0$. For $\ell \ge 1$, define the internal empirical feature covariance matrix by
\begin{equation}\label{eq:ch2-phi-ell-emp}
\Phi_\ell^{(N),\alpha\beta} := \frac{1}{N} \ip{\ns{\varphi}(h_\ell^\alpha)}{\ns{\varphi}(h_\ell^\beta)}.
\end{equation}
Let $\mathcal{F}_0$ be the trivial internal $\sigma$-field, and for $\ell \ge 1$ define the internal filtration (history)
\begin{equation}\label{eq:ch2-filtration}
\mathcal{F}_\ell := \ns{\sigma}\bigl(\{h_k^\alpha : \alpha \in [m],\, k \le \ell\}\bigr).
\end{equation}
Thus $\mathcal{F}_\ell$ records the internal hidden-neuron history through layer $\ell$; by internal independence across layers, $W_\ell, \dots, W_d$ are internally independent of $\mathcal{F}_\ell$. Conditioned on $\mathcal{F}_\ell$, the next-layer preactivations are internally Gaussian because $h_{\ell+1}^\alpha$ is an internal linear image of the Gaussian matrix $W_\ell$. More precisely, by Transfer of Proposition~\ref{prop:ch2-gauss-matvec}, for $\ell = 0, \dots, d-1$, the $m$-tuple $(h_{\ell+1}^\alpha)_{\alpha=1}^{m}$ (each in $\hR^N$) is internally conditionally Gaussian with covariance of the form
\begin{equation}\label{eq:ch2-hell-cond-law}
\bigl(h_{\ell+1}^\alpha\bigr)_{\alpha=1}^{m} \,\Big|\, \mathcal{F}_\ell \sim \hNcal\bigl(0, \Phi_\ell^{(N)} \otimes I_N\bigr).
\end{equation}

The key observation is that $\Phi_{\ell+1}^{(N)}$ is, up to infinitesimal fluctuations at unlimited width $N$, a deterministic function of $\Phi_\ell^{(N)}$. Indeed, by definition
\begin{equation}\label{eq:ch2-phi-ell1-sum}
\Phi_{\ell+1}^{(N),\alpha\beta}
= \frac{1}{N} \ip{\ns{\varphi}(h_{\ell+1}^\alpha)}{\ns{\varphi}(h_{\ell+1}^\beta)}
= \frac{1}{N} \sum_{j=1}^{N} \ns{\varphi}\bigl(h_{\ell+1,j}^\alpha\bigr) \, \ns{\varphi}\bigl(h_{\ell+1,j}^\beta\bigr).
\end{equation}
Conditioned on $\mathcal{F}_\ell$, the neuron slices $\{(h_{\ell+1,j}^\alpha)_{\alpha=1}^{m}\}_{j=1}^{N}$ are internally i.i.d.\ Gaussians in $\hR^m$ with covariance $\Phi_\ell^{(N)}$ (this is the meaning of~\eqref{eq:ch2-hell-cond-law}). Under the assumptions of Theorem~\ref{thm:ch2-deep-nngp}, for each fixed pair $(\alpha, \beta)$,
\begin{equation}\label{eq:ch2-phi-cond-conc}
\Phi_{\ell+1}^{(N),\alpha\beta} - \ns{\mathcal{C}}_\varphi\bigl(\Phi_\ell^{(N)}\bigr)^{\alpha\beta} \approxP 0,
\end{equation}
where the deterministic covariance map $\mathcal{C}_\varphi$ (with natural extension $\ns{\mathcal{C}}_\varphi$ by Transfer) is given entrywise by a Gaussian expectation:
\begin{equation}\label{eq:ch2-cphi-def}
\mathcal{C}_\varphi(\Phi)^{\alpha\beta} := \E\bigl[\varphi(g^\alpha)\varphi(g^\beta)\bigr], \qquad (g^1, \dots, g^m) \sim \mathcal{N}(0, \Phi).
\end{equation}
Define the deterministic standard forward kernels recursively by
\begin{equation}\label{eq:ch2-phi-rec}
\Phi_{\ell+1} := \mathcal{C}_\varphi(\Phi_\ell), \qquad \ell = 0, \dots, d-1.
\end{equation}
The theorem below makes precise that $\Phi_\ell^{(N)} \approxP \Phi_\ell$ for each standard $\ell \le d$.

\begin{remark}[Order of Limits]\label{rem:ch2-order-limits}
For deep networks, classical statements like ``$n \to \infty$'' implicitly hold depth fixed; if depth grows too, the order of limits matters, since taking $n \to \infty$ first and then $d \to \infty$ can differ from a joint limit. In nonstandard analysis, ``holding depth fixed at infinite width'' literally means that the depth $d \in \N$ is standard while the hidden width $N \in \hNinf$ is unlimited. If instead depth $D \in \hNinf$ is also unlimited, then sequential and joint limits become distinct hyperfinite regimes in $(\hN)^2$: $D/N \approx 0$ (width-dominant), $N/D \approx 0$ (depth-dominant), and $\st(D/N) = \bar{\tau} \in (0, \infty)$ (the proportional regime studied in Chapters~5 and~6).
\end{remark}

\begin{theorem}[Deep NNGP Recursion~\cite{lee2018deep,matthews2018gaussian}]\label{thm:ch2-deep-nngp}
Assume $\varphi : \R \to \R$ is continuous and has at most polynomial growth. Fix standard depth $d \in \N$, the standard number of inputs $m \in \N$, and the standard inputs $\{x^\alpha\}_{\alpha=1}^{m} \subset \R^{n_0}$ themselves. Then, at any unlimited hidden width $N \in \hNinf$,
\begin{equation}\label{eq:ch2-deep-nngp-conc}
\Phi_\ell^{(N)} \approxP \Phi_\ell \qquad \text{for each } \ell = 0, 1, \dots, d,
\end{equation}
where the deterministic matrices $\Phi_\ell$ are defined by~\eqref{eq:ch2-phi-rec}. Consequently, conditioned on the inputs, the output vector $\bigl(f(x^\alpha; \theta)\bigr)_{\alpha=1}^{m}$ belongs to $\Lim^m$ $\Loeb$-almost surely and its standard part $\st\bigl(f(x^1; \theta), \dots, f(x^m; \theta)\bigr)$ has exact Gaussian law $\mathcal{N}(0, \Phi_d)$ under $\Loeb$.
\end{theorem}

\begin{proofsketch}
The internal conditional law~\eqref{eq:ch2-hell-cond-law} implies that, given $\mathcal{F}_\ell$, the coordinates $\{(h_{\ell+1,j}^\alpha)_{\alpha=1}^{m}\}_{j=1}^{N}$ are internally i.i.d.\ copies of a centered Gaussian vector in $\hR^m$ with covariance $\Phi_\ell^{(N)}$. Therefore each entry of $\Phi_{\ell+1}^{(N)}$ is a hyperfinite empirical average of conditionally internally i.i.d.\ random variables,
\begin{equation}\label{eq:ch2-phi-ell1-avg}
\Phi_{\ell+1}^{(N),\alpha\beta} = \frac{1}{N} \sum_{j=1}^{N} \ns{\varphi}\bigl(h_{\ell+1,j}^\alpha\bigr)\ns{\varphi}\bigl(h_{\ell+1,j}^\beta\bigr),
\end{equation}
whose internal conditional expectation is $\ns{\mathcal{C}}_\varphi(\Phi_\ell^{(N)})^{\alpha\beta}$. Its internal conditional centered second moment is
\begin{equation}\label{eq:ch2-phi-cond-var}
\hE\left[ \left\lvert \Phi_{\ell+1}^{(N),\alpha\beta} - \ns{\mathcal{C}}_\varphi\bigl(\Phi_\ell^{(N)}\bigr)^{\alpha\beta} \right\rvert^2 \,\middle|\, \mathcal{F}_\ell \right]
= \frac{1}{N} \, \hVar_{g \sim \hNcal(0, \Phi_\ell^{(N)})}\!\Bigl(\ns{\varphi}(g^\alpha)\ns{\varphi}(g^\beta)\Bigr).
\end{equation}
We proceed by external induction over standard $\ell = 0, 1, \dots, d$, starting from $\Phi_0^{(N)} = \Phi_0$. Polynomial growth of $\varphi$ and the induction hypothesis $\Phi_\ell^{(N)} \approxP \Phi_\ell$ keep the internal conditional Gaussian variance on the right-hand side of~\eqref{eq:ch2-phi-cond-var} limited ($\in \Lim$) $\Loeb$-almost surely; dividing by $N \in \hNinf$ makes the conditional second moment infinitesimal ($\approx 0$), so internal Chebyshev's inequality yields~\eqref{eq:ch2-phi-cond-conc}. Continuity and polynomial growth also make $\mathcal{C}_\varphi$ continuous on bounded subsets of the standard positive semidefinite cone, so $\ns{\mathcal{C}}_\varphi$ is S-continuous at $\Phi_\ell$: $\Phi_\ell^{(N)} \approxP \Phi_\ell$ gives $\ns{\mathcal{C}}_\varphi(\Phi_\ell^{(N)}) \approxP \mathcal{C}_\varphi(\Phi_\ell) = \Phi_{\ell+1}$. Combining with~\eqref{eq:ch2-phi-cond-conc} completes the external induction over the finitely many standard layers $\ell = 0, 1, \dots, d$, giving $\Phi_\ell^{(N)} \approxP \Phi_\ell$ entrywise and hence as an $m \times m$ matrix (and $\Phi_d^{(N)} \approx \Phi_d$ $\Loeb$-almost surely by countable additivity of $\Loeb$). Finally, since $W_d$ is internally independent of $\mathcal{F}_d$, the output-layer identity~\eqref{eq:ch2-cond-gauss} gives, conditional on $\mathcal{F}_d$, the exact internal Gaussian output law $\hNcal(0, \Phi_d^{(N)})$; Loeb bounded convergence of the conditional characteristic function then yields the stated standard-part law $\mathcal{N}(0, \Phi_d)$.
\end{proofsketch}

\begin{remark}[Sequential versus Simultaneous Width Limits]\label{rem:ch2-seq-vs-sim}
The main derivation of Lee et al.~\cite{lee2018deep} takes the hidden-layer widths to infinity successively, while their Appendix~C gives an order-independent Gaussian-prior argument. Matthews et al.~\cite{matthews2018gaussian} give a rigorous simultaneous-width convergence theorem. In the internal Gaussian-weight, standard-depth, standard-finite-dataset setting used here, exact internal conditional Gaussianity at a single unlimited width $N \in \hNinf$ leads directly to the common-width covariance recursion above without iterated limits.
\end{remark}

\section{Neural Tangent Kernel and Linearized Training}\label{sec:ch2-ntk-general}

If we train the wide neural networks introduced above in the NTK regime, their output dynamics behave like those of kernel methods. The key object is the Neural Tangent Kernel (NTK), a kernel induced by the network's tangent features. At unlimited width $N \in \hNinf$, the empirical NTK becomes, up to infinitesimals, (i) deterministic and (ii) constant in time along gradient flow on standard time intervals, yielding closed-form linear standard-part training dynamics and, when the limiting kernel is positive definite, exponential convergence of the training loss. We then extend the construction from one-hidden-layer networks to finite-depth multilayer perceptrons (MLPs) and highlight the geometric interpretation as a preconditioned gradient flow in function space. We will reuse the forward-kernel notation and NNGP recursions from earlier sections.

The Neural Tangent Kernel was introduced by Jacot et al.~\cite{jacot2018neural}; Lee et al.~\cite{lee2019wide} made explicit the associated linearization of wide-network training. A closely related line of work provides quantitative convergence guarantees for overparameterized deep networks; see, for instance, Allen-Zhu et al.~\cite{allen2018convergence}, Du et al.~\cite{du2018gradient}, and Zou et al.~\cite{zou2018stochastic}. In these notes we emphasize the hyperfinite infinite-width picture and use these works primarily as pointers to the finite-width quantitative theory.

Let the standard training dataset (with $m, n_0 \in \N$ standard) be
\begin{equation}\label{eq:ch2-dataset}
\{(x^\alpha, y^\alpha)\}_{\alpha=1}^{m}, \qquad x^\alpha \in \R^{n_0}, \quad y^\alpha \in \R.
\end{equation}
Let $f(\cdot; \theta) : \hR^{n_0} \to \hR$ be an internal neural network with parameter vector $\theta \in \hR^{\dim(\theta)}$. We write $f_\alpha(\theta) := f(x^\alpha; \theta)$ and collect outputs into the internal vector
\begin{equation}\label{eq:ch2-f-y-vec}
f(\theta) := \bigl(f_1(\theta), \dots, f_m(\theta)\bigr)^\top \in \hR^m, \qquad y := (y^1, \dots, y^m)^\top \in \R^m.
\end{equation}
We will often use the internal empirical mean-squared error
\begin{equation}\label{eq:ch2-mse-loss}
L(\theta) := \frac{1}{2m} \sum_{\alpha=1}^{m} \bigl(f_\alpha(\theta) - y^\alpha\bigr)^2 = \frac{1}{2m} \norm{f(\theta) - y}_2^2.
\end{equation}
Define the internal residual vector $r(\theta) := f(\theta) - y \in \hR^m$.

Internal gradient descent with step size $\eta$ is
\begin{equation}\label{eq:ch2-gd}
\theta_{k+1} = \theta_k - \eta \, \nabla_\theta L(\theta_k),
\end{equation}
where $\nabla_\theta$ denotes the internal parameter gradient (obtained by Transfer). A convenient continuous-time counterpart is internal gradient flow, defined for $t \in \hR_{\ge 0}$ by Transfer of the standard ODE:
\begin{equation}\label{eq:ch2-gf}
\partial_t \theta(t) = -\nabla_\theta L(\theta(t)), \qquad \theta(0) = \theta_0,
\end{equation}
where $\partial_t \theta(t) = \frac{d}{dt}\theta(t)$ in the internal sense. In this chapter we analyze the induced dynamics of the network outputs $f(t) := f(\theta(t))$.

We now encode the network's parameter sensitivities through its internal tangent feature map. For each input $x$, define the tangent feature map at parameters $\theta$ by
\begin{equation}\label{eq:ch2-tangent-map}
\nabla_\theta f(x; \theta) \in \hR^{\dim(\theta)}.
\end{equation}
For a hyperfinite width-$N$ network, the (empirical) internal Neural Tangent Kernel at time $t \in \hR_{\ge 0}$ is the $m \times m$ matrix
\begin{equation}\label{eq:ch2-ntk-def}
K_t^{(N)} \in \hR^{m \times m}, \qquad K_t^{(N),\alpha\beta} := K_t^{(N)}(x^\alpha, x^\beta) := \ip{\nabla_\theta f_\alpha(\theta(t))}{\nabla_\theta f_\beta(\theta(t))}.
\end{equation}
By Transfer, $K_t^{(N)}$ is internally symmetric positive semidefinite (PSD).

By the Transfer Principle, a chain-rule computation yields a closed-form internal ODE for $r(t)$ in terms of $K_t^{(N)}$.

\begin{proposition}[Residual ODE]\label{prop:ch2-residual-ode}
Under internal gradient flow for the squared loss, the residual vector $r(t) = f(t) - y \in \hR^m$ satisfies, for all $t \in \hR_{\ge 0}$,
\begin{equation}\label{eq:ch2-residual-ode}
\partial_t r(t) = -\frac{1}{m} K_t^{(N)} r(t).
\end{equation}
\end{proposition}

\begin{proof}
By Transfer of the chain rule, for each $\alpha \in [m]$,
\begin{equation}\label{eq:ch2-res-ode-step1}
\partial_t f_\alpha(t) = \ip{\nabla_\theta f_\alpha(\theta(t))}{\partial_t \theta(t)} = -\ip{\nabla_\theta f_\alpha(\theta(t))}{\nabla_\theta L(\theta(t))}.
\end{equation}
For the squared loss, $\nabla_\theta L(\theta) = \frac{1}{m} \sum_{\beta=1}^{m} (f_\beta(\theta) - y^\beta)\nabla_\theta f_\beta(\theta)$. Substituting gives
\begin{equation}\label{eq:ch2-res-ode-step2}
\partial_t f_\alpha(t) = -\frac{1}{m} \sum_{\beta=1}^{m} \ip{\nabla_\theta f_\alpha(\theta(t))}{\nabla_\theta f_\beta(\theta(t))}\bigl(f_\beta(t) - y^\beta\bigr) = -\frac{1}{m} \sum_{\beta=1}^{m} K_t^{(N),\alpha\beta} r_\beta(t).
\end{equation}
In vector form this is~\eqref{eq:ch2-residual-ode}.
\end{proof}

\begin{remark}\label{rem:ch2-residual-mixing}
Equation~\eqref{eq:ch2-residual-ode} shows that each residual coordinate $r_\alpha(t)$ mixes information from all residuals through the internal kernel matrix $K_t^{(N)}$. If $K_t^{(N)}$ were replaced by a fixed standard kernel matrix $K$, then training would become a linear ODE in $r$ with an explicit solution via a matrix exponential. The main content of the NTK theory is that at unlimited width $N \in \hNinf$, $K_t^{(N)} \approx K$ uniformly on standard time intervals, making this replacement exact on standard parts.
\end{remark}

This framework easily extends to general loss functions. For a differentiable pointwise loss $\ell(f, y)$ (extended internally via $\ns{\ell}$), define the internal empirical risk
\begin{equation}\label{eq:ch2-general-risk}
L(\theta) := \frac{1}{m} \sum_{\alpha=1}^{m} \ns{\ell}\bigl(f_\alpha(\theta), y^\alpha\bigr).
\end{equation}
Along the internal training trajectory, define $r_\alpha(t) := \partial_f \ns{\ell}(f_\alpha(t), y^\alpha)$ and collect these entries in $r(t) \in \hR^m$. For squared loss, this agrees with the residual $f(t) - y$ defined above.

\begin{proposition}[Risk Dissipation Identity]\label{prop:ch2-risk-dissipation}
Along internal gradient flow for $L$,
\begin{equation}\label{eq:ch2-risk-dissipation}
\frac{d}{dt} L(\theta(t)) = -\frac{1}{m^2} r(t)^\top K_t^{(N)} r(t).
\end{equation}
\end{proposition}

\begin{proof}
By Transfer, the chain rule gives
\begin{equation}\label{eq:ch2-risk-diss-1}
\frac{d}{dt} L(\theta(t)) = \frac{1}{m} \sum_{\alpha=1}^{m} r_\alpha(t) \, \partial_t f_\alpha(t).
\end{equation}
Moreover,
\begin{equation}\label{eq:ch2-risk-diss-2}
\nabla_\theta L(\theta(t)) = \frac{1}{m} \sum_{\beta=1}^{m} r_\beta(t) \, \nabla_\theta f_\beta(\theta(t)).
\end{equation}
Substituting $\partial_t f_\alpha(t) = -\ip{\nabla_\theta f_\alpha(\theta(t))}{\nabla_\theta L(\theta(t))}$ and using~\eqref{eq:ch2-ntk-def} yields~\eqref{eq:ch2-risk-dissipation}.
\end{proof}

\section{The One-Layer Case}\label{sec:ch2-one-layer-ntk}

The empirical NTK $K_t^{(N)}$ is random at initialization (because the parameters are random) and evolves during training (because $\theta(t)$ changes). At unlimited width $N \in \hNinf$, the NTK regime exhibits two related features:
\begin{enumerate}[label=(\roman*)]
\item \emph{Determinism:} for every standard training time $t \ge 0$, the empirical kernel $K_t^{(N)}(x, x')$ concentrates infinitesimally around a deterministic standard kernel $K(x, x')$.
\item \emph{Stationarity:} the limiting standard kernel $K$ does not depend on $t$.
\end{enumerate}
We start with the simplest case to illustrate the mechanism.

\begin{definition}[One-Hidden-Layer MLP (NTK Parameterization)]\label{def:ch2-1layer-ntk}
Fix an unlimited width $N \in \hNinf$ and standard input dimension $n_0 \in \N$. Let
\begin{equation}\label{eq:ch2-1layer-ntk-model}
f(x; \theta) = \frac{1}{\sqrt{N}} a^\top \ns{\varphi}\!\left(\frac{Wx}{\sqrt{n_0}}\right), \qquad W \in \hR^{N \times n_0}, \quad a \in \hR^N,
\end{equation}
where $\varphi : \R \to \R$ (extended via $\ns{\varphi}$) is applied elementwise to vectors. We take internally i.i.d.\ Gaussian initialization $W_{ij} \sim \hNcal(0, 1)$ and $a_i \sim \hNcal(0, 1)$ independently on $(\Omega, \mathcal{A}, \hP)$.
\end{definition}

\begin{remark}[Bias Parameters]\label{rem:ch2-biases}
For notational simplicity we omit biases throughout this set of notes. Including bias terms does not change any of the core ideas (infinitesimal concentration of kernels, near-stationarity of the tangent kernel, etc.); it mainly makes the algebra longer and more tedious in an uninformative way.
\end{remark}

\begin{theorem}[One-Hidden-Layer NTK Limit (Informal)]\label{thm:ch2-1layer-ntk-thm}
Consider the width-$N$ network ($N \in \hNinf$) and NTK parameterization in~\eqref{eq:ch2-1layer-ntk-model}, with the internal Gaussian initialization above, trained by internal gradient flow for squared loss on a fixed standard finite dataset. Write $K_s^{(N)}$ for its empirical NTK at internal time $s \in \hR_{\ge 0}$. For sufficiently nice $\varphi$, the NTK regime~\cite{jacot2018neural,lee2019wide} gives a deterministic standard symmetric positive semidefinite matrix $K \in \R^{m \times m}$, independent of $t$, such that for every standard $t \ge 0$ (identifying $t$ with $\ns{t} \in \hR_{\ge 0}$),
\begin{equation}\label{eq:ch2-1layer-ntk-conc}
\sup_{0 \le s \le \ns{t}} \norm{K_s^{(N)} - K} \approxP 0, \qquad
\st\bigl(K_s^{(N)}\bigr) = K \quad (0 \le s \le \ns{t} \text{ $\Loeb$-a.s.}).
\end{equation}
The standard-part residual dynamics satisfy
\begin{equation}\label{eq:ch2-1layer-ntk-res-ode}
\partial_t \st\bigl(r(t)\bigr) = -\frac{1}{m} K \, \st\bigl(r(t)\bigr), \qquad \text{so} \qquad \st\bigl(r(t)\bigr) = \exp\!\left(-\frac{t}{m} K\right)\st\bigl(r(0)\bigr).
\end{equation}
\end{theorem}

\begin{proofsketch}
Fix a standard finite dataset $\{x^\alpha\}_{\alpha=1}^{m}$.

\emph{Step 1: Determinism.} Writing $w_i^\top$ for the $i$-th row of $W$, for any internal parameters $\theta$ (and in particular at $\theta_0$), the neuronwise internal gradients for a single datum $x^\alpha$ are
\begin{equation}\label{eq:ch2-1layer-neuron-grad}
\frac{\partial f_\alpha}{\partial a_i} = \frac{1}{\sqrt{N}} \ns{\varphi}(h_i^\alpha), \qquad
\frac{\partial f_\alpha}{\partial w_i} = \frac{1}{\sqrt{N n_0}} a_i \, \ns{\varphi}'(h_i^\alpha) x^\alpha, \qquad
h_i^\alpha := \frac{1}{\sqrt{n_0}} w_i^\top x^\alpha.
\end{equation}
Plugging~\eqref{eq:ch2-1layer-neuron-grad} into $K_0^{(N),\alpha\beta} = \ip{\nabla_\theta f_\alpha(\theta_0)}{\nabla_\theta f_\beta(\theta_0)}$ yields the exact internal decomposition
\begin{equation}\label{eq:ch2-1layer-kappa-sum}
K_0^{(N),\alpha\beta} = \frac{1}{N} \sum_{i=1}^{N} \kappa_i^{\alpha\beta}, \qquad
\kappa_i^{\alpha\beta} := \ns{\varphi}(h_i^\alpha)\ns{\varphi}(h_i^\beta) + a_i^2 \, \ns{\varphi}'(h_i^\alpha)\ns{\varphi}'(h_i^\beta) \frac{1}{n_0} \ip{x^\alpha}{x^\beta},
\end{equation}
where $\kappa_i^{\alpha\beta}$ is evaluated at initialization $\theta = \theta_0$. Because $(a_i, w_i)$ are internally i.i.d.\ at Gaussian initialization, for each fixed pair $(\alpha, \beta)$ the internal random variables $\{\kappa_i^{\alpha\beta}\}_{i=1}^{N}$ are internally i.i.d., and hence, provided the standard first moment $\E\abs{\kappa_1^{\alpha\beta}} < \infty$, the hyperfinite law of large numbers gives $K_0^{(N),\alpha\beta} \approxP K^{\alpha\beta}$ (in fact $\Loeb$-almost surely) for each fixed $(\alpha, \beta)$. Since $m \in \N$ is standard, this entrywise infinitesimal concentration implies $\norm{K_0^{(N)} - K} \approxP 0$ in any standard matrix norm. Moreover, $K$ is positive semidefinite because each $K_0^{(N)}$ is internally positive semidefinite and taking standard parts preserves $u^\top K_0^{(N)} u \ge 0$ for all $u \in \R^m$.

The neuronwise form in~\eqref{eq:ch2-1layer-kappa-sum} is convenient for the law-of-large-numbers argument and will also be used in the stationarity estimate below. To display separately the contributions of the two parameter blocks, write $h^\alpha = (h_1^\alpha, \dots, h_N^\alpha)^\top \in \hR^N$. Then the same identity becomes
\begin{equation}\label{eq:ch2-1layer-block-ntk}
\begin{aligned}
K_0^{(N),\alpha\beta} &= \frac{1}{N} \ip{\ns{\varphi}(h^\alpha)}{\ns{\varphi}(h^\beta)} + \Phi_0^{\alpha\beta} \frac{1}{N} \ip{a \odot \ns{\varphi}'(h^\alpha)}{a \odot \ns{\varphi}'(h^\beta)} \\
&\approxP K^{\alpha\beta} = \Phi_1^{\alpha\beta} + \Phi_0^{\alpha\beta} \Phi_1^{\prime,\alpha\beta}.
\end{aligned}
\end{equation}
Here $\Phi_1^{\prime,\alpha\beta} = \E[\varphi'(g^\alpha)\varphi'(g^\beta)]$ for a centered standard Gaussian pair with covariance $[\Phi_0^{\gamma\delta}]_{\gamma,\delta \in \{\alpha, \beta\}}$.

\emph{Step 2: Stationarity on a standard time interval.} Fix a standard $t \ge 0$. For the one-hidden-layer model, internal gradient flow gives, for $0 \le s \le \ns{t}$,
\begin{equation}\label{eq:ch2-1layer-gf-coords}
\partial_s a_i(s) = -\frac{1}{m\sqrt{N}} \sum_{\gamma=1}^{m} r_\gamma(s) \ns{\varphi}\bigl(h_i^\gamma(s)\bigr), \qquad
\partial_s w_i(s) = -\frac{a_i(s)}{m\sqrt{N n_0}} \sum_{\gamma=1}^{m} r_\gamma(s) \ns{\varphi}'\bigl(h_i^\gamma(s)\bigr) x^\gamma.
\end{equation}
Along internal gradient flow, by Transfer, the decrease in loss $L(\theta_0) - L(\theta(s)) \le L(\theta_0)$ equals the internal time integral $\int_0^s \norm{\partial_u \theta(u)}_2^2\, du$ of the squared parameter velocity. Internal Cauchy--Schwarz therefore gives the empirical parameter displacement bound
\begin{equation}\label{eq:ch2-1layer-disp-bound}
\sup_{0 \le s \le \ns{t}} \frac{1}{N} \sum_{i=1}^{N} \Bigl[(a_i(s) - a_i(0))^2 + \norm{w_i(s) - w_i(0)}_2^2\Bigr] \le \frac{t L(\theta_0)}{N} \approx 0,
\end{equation}
an infinitesimal of order $N^{-1}$ $\Loeb$-almost surely because $L(\theta_0) \in \Lim$ $\Loeb$-almost surely under the Gaussian initialization on the fixed standard dataset. Let $\kappa_i^{\alpha\beta}(s)$ denote the expression in~\eqref{eq:ch2-1layer-kappa-sum} evaluated along the internal flow at $s \in [0, \ns{t}]$. An internal first-order Taylor expansion at initialization gives the next-order correction
\begin{equation}\label{eq:ch2-1layer-taylor}
\begin{aligned}
K_s^{(N),\alpha\beta} - K_0^{(N),\alpha\beta}
&= \frac{1}{N} \sum_{i=1}^{N} \bigl(\kappa_i^{\alpha\beta}(s) - \kappa_i^{\alpha\beta}(0)\bigr) \\
&= \frac{1}{N} \sum_{i=1}^{N} \Bigl[ \partial_{a_i} \kappa_i^{\alpha\beta}(0)\bigl(a_i(s) - a_i(0)\bigr) \\
&\qquad\qquad + \ip{\nabla_{w_i}\kappa_i^{\alpha\beta}(0)}{w_i(s) - w_i(0)} \Bigr] + R_s^{(N),\alpha\beta},
\end{aligned}
\end{equation}
with internal remainder satisfying $\sup_{0 \le s \le \ns{t}} N^{1/2}\abs{R_s^{(N),\alpha\beta}} \approxP 0$. The derivatives in the leading term are
\begin{align}
\frac{\partial \kappa_i^{\alpha\beta}}{\partial a_i} &= 2 a_i \ns{\varphi}'(h_i^\alpha)\ns{\varphi}'(h_i^\beta) \frac{\ip{x^\alpha}{x^\beta}}{n_0}, \label{eq:ch2-1layer-dkappa-da} \\
\nabla_{w_i} \kappa_i^{\alpha\beta} &= \frac{x^\alpha}{\sqrt{n_0}} \left( \ns{\varphi}'(h_i^\alpha)\ns{\varphi}(h_i^\beta) + a_i^2 \ns{\varphi}''(h_i^\alpha)\ns{\varphi}'(h_i^\beta) \frac{\ip{x^\alpha}{x^\beta}}{n_0} \right) \nonumber \\
&\quad + \frac{x^\beta}{\sqrt{n_0}} \left( \ns{\varphi}(h_i^\alpha)\ns{\varphi}'(h_i^\beta) + a_i^2 \ns{\varphi}'(h_i^\alpha)\ns{\varphi}''(h_i^\beta) \frac{\ip{x^\alpha}{x^\beta}}{n_0} \right). \label{eq:ch2-1layer-dkappa-dw}
\end{align}
For sufficiently nice $\varphi$, these derivatives have limited empirical second moments ($\in \Lim$) at initialization $\Loeb$-almost surely. Internal Cauchy--Schwarz and~\eqref{eq:ch2-1layer-disp-bound} therefore bound the averaged linear Taylor term by an infinitesimal of scale $N^{-1/2}$ ($\in N^{-1/2}\Lim$), uniformly on $[0, \ns{t}]$. The smoothness and moment assumptions also control the internal Taylor remainder $R_s^{(N),\alpha\beta}$ uniformly on $[0, \ns{t}]$, giving the same $N^{-1/2}$-scale infinitesimal bound for the full correction $K_s^{(N),\alpha\beta} - K_0^{(N),\alpha\beta}$ in~\eqref{eq:ch2-1layer-taylor}. Combining this with Step~1 gives~\eqref{eq:ch2-1layer-ntk-conc}, so $\st(K_s^{(N)}) = K$ uniformly for $0 \le s \le \ns{t}$. Substituting $K_s^{(N)} \approx K$ into the internal time-integral form of~\eqref{eq:ch2-residual-ode} via internal Gr\"onwall's inequality, together with the NNGP shadow for the initial residual $r(0)$, gives $r(\ns{t}) \approx \ns{\exp}(-\frac{t}{m}K)r(0)$ and hence the standard-part residual dynamics~\eqref{eq:ch2-1layer-ntk-res-ode}.
\end{proofsketch}

\begin{remark}[Scope of the Proof Sketch]\label{rem:ch2-scope-sketch}
The calculation above displays the one-hidden-layer mechanism. For multiple hidden layers, an additional layerwise argument is required; we omit those details here. Nonsmooth activations such as ReLU require a different argument.
\end{remark}

We now study the limiting fixed-kernel dynamics, in which the time-dependent empirical kernel in~\eqref{eq:ch2-residual-ode} is replaced by a fixed standard positive-definite matrix $K$. The resulting linear ODE has an explicit solution and converges exponentially.

\begin{proposition}[Exponential Convergence of the Fixed-Kernel Dynamics]\label{prop:ch2-fixed-kernel-exp}
Let $K \in \R^{m \times m}$ be symmetric positive definite, and suppose $\partial_t r(t) = -m^{-1} K r(t)$. Then
\begin{equation}\label{eq:ch2-fixed-kernel-sol}
r(t) = \exp\!\left(-\frac{t}{m} K\right) r(0),
\end{equation}
and the squared loss $L(t) := (2m)^{-1}\norm{r(t)}_2^2$ satisfies
\begin{equation}\label{eq:ch2-fixed-kernel-loss}
L(t) \le \exp\!\left(-\frac{2\lambda_{\min}(K)}{m} t\right) L(0).
\end{equation}
Thus the loss converges to zero exponentially.
\end{proposition}

\begin{proof}
By the spectral theorem, $K = P \diag(\lambda_1, \dots, \lambda_m) P^\top$ for an orthogonal matrix $P$ and eigenvalues $\lambda_i > 0$. The matrix exponential in~\eqref{eq:ch2-fixed-kernel-sol} solves the constant-coefficient residual ODE. Moreover,
\begin{equation}\label{eq:ch2-fixed-kernel-norm}
\norm{r(t)}_2^2 = \sum_{i=1}^{m} e^{-2\lambda_i t / m} \bigl((P^\top r(0))_i\bigr)^2 \le e^{-2\lambda_{\min}(K) t / m} \norm{r(0)}_2^2.
\end{equation}
Dividing by $2m$ proves~\eqref{eq:ch2-fixed-kernel-loss}.
\end{proof}

\begin{remark}[When Is $\lambda_{\min}(K) > 0$?]\label{rem:ch2-lambda-min}
For a fixed standard dataset, positive definiteness is a property of the deterministic standard kernel $K$. If $\lambda_{\min}(K) > 0$, then $K_0^{(N)} \approxP K$ and continuity of $\lambda_{\min}$ imply $\lambda_{\min}(K_0^{(N)}) \approxP \lambda_{\min}(K) > 0$, so
\begin{equation}\label{eq:ch2-k0-pd}
\hP\bigl(K_0^{(N)} \text{ is positive definite}\bigr) \approx 1, \qquad
\Loeb\bigl(K_0^{(N)} \text{ is positive definite}\bigr) = 1.
\end{equation}
When the data are random, positivity and conditioning of $K$ can themselves be studied probabilistically. See, for example, Nguyen et al.~\cite{nguyen2021tight} for quantitative lower bounds on the smallest eigenvalue of empirical NTK matrices for deep ReLU networks under standard assumptions.
\end{remark}

\begin{remark}[Linearization and the Kernel Regression Viewpoint]\label{rem:ch2-linearization}
The NTK regime can be viewed directly as an infinitesimal linearization of the network around initialization. Consider internal gradient descent with standard step size $\eta$: $\theta_{k+1} = \theta_k - \eta \nabla_\theta L(\theta_k)$. At unlimited width $N \in \hNinf$, the hyperfinite NTK width counting gives
\begin{equation}\label{eq:ch2-ntk-step-taylor}
f_\alpha(\theta_{k+1}) = f_\alpha(\theta_k) + \ip{\nabla_\theta f_\alpha(\theta_k)}{\theta_{k+1} - \theta_k} + \underbrace{\varepsilon_k^\alpha}_{\in N^{-1/2}\Lim\text{, nonlinear correction}}.
\end{equation}
Here the $\varepsilon_k^\alpha \in N^{-1/2}\Lim \subset \Inf$ term summarizes the leading nonlinear correction, an infinitesimal of scale $N^{-1/2}$. Since this correction is infinitesimal and the tangent features remain infinitesimally frozen along standard-time training, we obtain the hyperreal linearization
\begin{equation}\label{eq:ch2-ntk-linearized}
f_\alpha(\theta) \approx f_\alpha(\theta_0) + \ip{\nabla_\theta f_\alpha(\theta_0)}{\theta - \theta_0}.
\end{equation}
Hence training becomes, up to infinitesimals, linear regression in the initial tangent features $\nabla_\theta f(x^\alpha; \theta_0)$. Rigorous NTK results make this linearization precise along the training trajectory~\cite{lee2019wide}.

A complementary geometric viewpoint is already visible in~\eqref{eq:ch2-residual-ode}. When the standard kernel $K$ is positive definite, equip the standard output space $\R^m$ with the inner product
\begin{equation}\label{eq:ch2-k-inv-metric}
\ip{u}{v}_{K^{-1}} := u^\top K^{-1} v.
\end{equation}
With respect to this inner product, the fixed-kernel standard-part dynamics become
\begin{equation}\label{eq:ch2-riemannian-gf}
\partial_t f = -\frac{1}{m} K(f - y) = -\grad_{K^{-1}}\!\left(\frac{1}{2m} \norm{f - y}_2^2\right).
\end{equation}
On the fixed dataset, the space of predictor values is simply the linear vector space $\R^m$. Nevertheless, $K^{-1}$ is precisely a constant, and hence flat, Riemannian metric on this space; equivalently, $K$ acts as its inverse metric or as a preconditioner.
\end{remark}

\section{Extension to Finite-Depth MLPs}\label{sec:ch2-deep-ntk}

We now describe the exact internal layerwise structure of the NTK for a standard-depth-$d$ fully-connected network of unlimited width $N \in \hNinf$ and then compute its deterministic standard shadow at Gaussian initialization.

The finite-depth NTK is built from three internal layerwise quantities: the empirical forward kernels $\Phi_\ell^{(N)}$, the empirical derivative kernels $\Phi_\ell^{\prime(N)}$, and the backward covariances $G_\ell^{(N)}$. At Gaussian initialization, the deterministic forward- and derivative-kernel shadows feed into a backward recursion for the shadows of the backward covariances. Substituting these shadows into the exact internal layerwise decomposition in Proposition~\ref{prop:ch2-layerwise-ntk} then yields the limiting standard NTK. Throughout this section, the superscript $(N)$ denotes hyperfinite width $N \in \hNinf$, and common training-time arguments are suppressed.

To analyze the finite-depth NTK, we use the following network notation. Let $N \in \hNinf$ be the unlimited hidden width and $n_0 \in \N$ the standard input dimension, with standard depth $d \in \N$. For $\ell = 0, 1, \dots, d$, let $W_\ell$ be internal weight matrices with internally i.i.d.\ $\hNcal(0, 1)$ entries, where $W_0 \in \hR^{N \times n_0}$, $W_\ell \in \hR^{N \times N}$ for $1 \le \ell \le d-1$, and $W_d \in \hR^{1 \times N}$ (scalar output). Define for standard input $x^\alpha$ the internal pre-activations
\begin{equation}\label{eq:ch2-deep-ntk-preact}
h_1^\alpha = \frac{1}{\sqrt{n_0}} W_0 x^\alpha, \qquad h_{\ell+1}^\alpha = \frac{1}{\sqrt{N}} W_\ell \ns{\varphi}(h_\ell^\alpha) \quad (1 \le \ell \le d-1),
\end{equation}
and output
\begin{equation}\label{eq:ch2-deep-ntk-out}
f_\alpha = f(x^\alpha; \theta) = \frac{1}{\sqrt{N}} W_d \ns{\varphi}(h_d^\alpha).
\end{equation}

We first record the forward and derivative kernels, reusing the forward-kernel notation and deterministic recursion from Section~\ref{sec:ch2-deep-nngp}. Define the standard input (zeroth-layer) kernel
\begin{equation}\label{eq:ch2-deep-ntk-phi0}
\Phi_0^{\alpha\beta} := \frac{1}{n_0} \ip{x^\alpha}{x^\beta}.
\end{equation}
Recall that $\Phi_0^{(N)} = \Phi_0$. For $\ell \ge 1$, define the internal empirical forward and derivative kernels
\begin{equation}\label{eq:ch2-deep-ntk-emp-kernels}
\Phi_\ell^{(N),\alpha\beta} := \frac{1}{N} \ip{\ns{\varphi}(h_\ell^\alpha)}{\ns{\varphi}(h_\ell^\beta)}, \qquad
\Phi_\ell^{\prime(N),\alpha\beta} := \frac{1}{N} \ip{\ns{\varphi}'(h_\ell^\alpha)}{\ns{\varphi}'(h_\ell^\beta)}.
\end{equation}
Theorem~\ref{thm:ch2-deep-nngp} gives $\Phi_\ell^{(N)} \approxP \Phi_\ell$. Suppose additionally that $\varphi'$ exists almost everywhere, its Gaussian-product second moments are uniformly bounded for covariance matrices in a standard neighborhood of the deterministic recursion, and the corresponding Gaussian expectations are continuous (hence S-continuous) there. Then the same internal conditional second-moment argument gives $\Phi_\ell^{\prime(N)} \approxP \Phi_\ell'$.

For $\ell \ge 0$, let $(h^\alpha, h^\beta)$ be the centered standard Gaussian pair whose covariance is determined by $\Phi_\ell$. The deterministic derivative-kernel shadow obeys $\Phi_{\ell+1}^{\prime,\alpha\beta} = \E[\varphi'(h^\alpha)\varphi'(h^\beta)]$. We write this compactly as $\Phi_{\ell+1}' = \mathcal{C}_{\varphi'}(\Phi_\ell)$, where $\mathcal{C}_{\varphi'}$ denotes the companion Gaussian expectation rather than the derivative of $\mathcal{C}_\varphi$.

For ReLU, the Cho--Saul formulas in Proposition~\ref{prop:ch2-cho-saul} make the forward and derivative-kernel recursions explicit. For fixed $\ell, \alpha, \beta$ with $\Phi_\ell^{\alpha\alpha}, \Phi_\ell^{\beta\beta} > 0$, write $\rho := \Phi_\ell^{\alpha\beta} / \sqrt{\Phi_\ell^{\alpha\alpha}\Phi_\ell^{\beta\beta}} \in [-1, 1]$. Then
\begin{equation}\label{eq:ch2-relu-rec}
\Phi_{\ell+1}^{\alpha\beta} = \frac{\sqrt{\Phi_\ell^{\alpha\alpha}\Phi_\ell^{\beta\beta}}}{2\pi}\left( \sqrt{1 - \rho^2} + \rho\left(\frac{\pi}{2} + \arcsin(\rho)\right) \right), \qquad
\Phi_{\ell+1}^{\prime,\alpha\beta} = \frac{1}{4} + \frac{1}{2\pi}\arcsin(\rho).
\end{equation}
Thus the same correlation determines both the next forward kernel and the derivative factor that enters the backward recursion.

We next introduce the internal backpropagated sensitivities whose covariances complete the NTK decomposition. Fix an internal training time $t \in \hR_{\ge 0}$ and evaluate all forward and backward quantities below at $\theta(t)$, suppressing their common time argument except on
\begin{equation}\label{eq:ch2-deep-ntk-def}
K_t^{(N),\alpha\beta} := \ip{\nabla_\theta f_\alpha(\theta(t))}{\nabla_\theta f_\beta(\theta(t))}.
\end{equation}
The additional state variables in the decomposition are the internal empirical backward covariances $G_\ell^{(N),\alpha\beta}$, which encode inner products of backpropagated sensitivities. At Gaussian initialization, the empirical forward kernels $\Phi_\ell^{(N),\alpha\beta}$ concentrate infinitesimally around the deterministic NNGP kernels $\Phi_\ell^{\alpha\beta}$ defined earlier. Define the internal backward sensitivity vectors
\begin{equation}\label{eq:ch2-g-def}
g_\ell^\alpha := \sqrt{N} \frac{\partial f_\alpha}{\partial h_\ell^\alpha} \in \hR^N, \qquad 1 \le \ell \le d.
\end{equation}
By Transfer, a direct backpropagation calculation gives the exact internal top-layer identity and hidden-layer recursion
\begin{align}
g_d^\alpha &= W_d^\top \odot \ns{\varphi}'(h_d^\alpha), \label{eq:ch2-gd-top} \\
g_\ell^\alpha &= \frac{1}{\sqrt{N}} \diag\bigl(\ns{\varphi}'(h_\ell^\alpha)\bigr) W_\ell^\top g_{\ell+1}^\alpha, \qquad \ell = d-1, \dots, 1. \label{eq:ch2-gell-rec}
\end{align}
Define the internal backward covariance matrices
\begin{equation}\label{eq:ch2-G-emp}
G_\ell^{(N),\alpha\beta} := \frac{1}{N} \ip{g_\ell^\alpha}{g_\ell^\beta}, \qquad \ell = 1, \dots, d, \qquad G_{d+1}^{(N),\alpha\beta} := 1.
\end{equation}

\begin{proposition}[Layerwise NTK Decomposition]\label{prop:ch2-layerwise-ntk}
For the depth-$d$ MLP defined above (with no biases), at every internal training time $t \in \hR_{\ge 0}$ the empirical NTK on the dataset $\{x^\alpha\}_{\alpha=1}^{m}$ decomposes exactly as
\begin{equation}\label{eq:ch2-layerwise-ntk}
K_t^{(N),\alpha\beta} = \sum_{\ell=0}^{d} G_{\ell+1}^{(N),\alpha\beta} \Phi_\ell^{(N),\alpha\beta},
\end{equation}
where $\Phi_0^{(N)} = \Phi_0$, and $G_{d+1}^{(N),\alpha\beta} = 1$ corresponds to the top (linear) output layer.

In matrix form, viewing $\Phi_\ell^{(N)}, G_{\ell+1}^{(N)}, K_t^{(N)} \in \hR^{m \times m}$,~\eqref{eq:ch2-layerwise-ntk} reads
\begin{equation}\label{eq:ch2-layerwise-ntk-hadamard}
K_t^{(N)} = \sum_{\ell=0}^{d} G_{\ell+1}^{(N)} \odot \Phi_\ell^{(N)},
\end{equation}
where $\odot$ denotes the Hadamard (entrywise) product.
\end{proposition}

\begin{proof}
By Transfer, the internal NTK is an inner product of parameter gradients: $K_t^{(N),\alpha\beta} = \ip{\nabla_\theta f_\alpha(\theta(t))}{\nabla_\theta f_\beta(\theta(t))}$. Split $\theta$ layer-by-layer and sum the resulting contributions. The definition of $g_\ell^\alpha$ and the network scaling give the boundary-aware internal identities
\begin{equation}\label{eq:ch2-layer-grads}
\begin{aligned}
\frac{\partial f_\alpha}{\partial W_0} &= \frac{1}{\sqrt{N n_0}} g_1^\alpha (x^\alpha)^\top, &
\frac{\partial f_\alpha}{\partial W_\ell} &= \frac{1}{N} g_{\ell+1}^\alpha \bigl(\ns{\varphi}(h_\ell^\alpha)\bigr)^\top \quad (1 \le \ell \le d-1), \\
\frac{\partial f_\alpha}{\partial W_d} &= \frac{1}{\sqrt{N}} \bigl(\ns{\varphi}(h_d^\alpha)\bigr)^\top. & &
\end{aligned}
\end{equation}
Their internal Frobenius inner products are therefore
\begin{equation}\label{eq:ch2-layer-frob}
\begin{aligned}
\ip{\frac{\partial f_\alpha}{\partial W_0}}{\frac{\partial f_\beta}{\partial W_0}}_F &= G_1^{(N),\alpha\beta} \Phi_0^{\alpha\beta}, &
\ip{\frac{\partial f_\alpha}{\partial W_\ell}}{\frac{\partial f_\beta}{\partial W_\ell}}_F &= G_{\ell+1}^{(N),\alpha\beta} \Phi_\ell^{(N),\alpha\beta} \quad (1 \le \ell \le d-1), \\
\ip{\frac{\partial f_\alpha}{\partial W_d}}{\frac{\partial f_\beta}{\partial W_d}}_F &= \Phi_d^{(N),\alpha\beta}. & &
\end{aligned}
\end{equation}
Since $\Phi_0^{(N)} = \Phi_0$ and $G_{d+1}^{(N),\alpha\beta} = 1$, summing these contributions yields~\eqref{eq:ch2-layerwise-ntk}.
\end{proof}

To compute the standard shadow of this decomposition at unlimited width $N \in \hNinf$, we must handle the internal dependence created by reusing the same weights in the forward and backward passes. Conditioning on the forward activations breaks the internal i.i.d.\ structure of those weights; the transferred Gaussian conditioning identity below isolates a usable internally independent component.

Write $\varphi_\ell^\alpha := \ns{\varphi}(h_\ell^\alpha)$ and collect the dataset postactivations into
\begin{equation}\label{eq:ch2-postact-mat}
\varphi_\ell := \begin{bmatrix} \varphi_\ell^1 & \cdots & \varphi_\ell^m \end{bmatrix} \in \hR^{N \times m}.
\end{equation}
Since $h_{\ell+1}^\alpha = \frac{1}{\sqrt{N}} W_\ell \varphi_\ell^\alpha$, internal conditioning on the next-layer preactivations is (up to the internal scaling $1/\sqrt{N}$) the same as conditioning on $W_\ell \varphi_\ell$.

\begin{lemma}[Gaussian Conditioning (One-Sided)]\label{lem:ch2-gauss-cond}
Let $W \in \hR^{N \times P}$ ($N, P \in \hN$) have internally i.i.d.\ $\hNcal(0, 1)$ entries and let $\varphi \in \hR^{P \times q}$ ($q \in \hN$) be internally deterministic. Let $P_\varphi := \varphi(\varphi^\top \varphi)^\dagger \varphi^\top$ be the internal orthogonal projector onto $\col(\varphi)$, and $P_\varphi^\perp := I_P - P_\varphi$. Then, conditioning internally on the internal $\sigma$-field $\ns{\sigma}(W\varphi)$,
\begin{equation}\label{eq:ch2-gauss-cond}
W \,\Big|\, \ns{\sigma}(W\varphi) \eqd W P_\varphi + \widetilde{W} P_\varphi^\perp,
\end{equation}
in internal law under $\hP$, where $\widetilde{W}$ is an internally independent copy of $W$.
\end{lemma}

\begin{proof}
See Appendix~A.5.
\end{proof}

\begin{proposition}[Backward Covariance Recursion]\label{prop:ch2-backward-rec}
Fix standard $m, d, n_0 \in \N$ and standard inputs $\{x^\alpha\}_{\alpha=1}^{m}$, and evaluate all quantities at internal Gaussian initialization with unlimited width $N \in \hNinf$. For sufficiently nice $\varphi$, there are deterministic standard matrices $G_1, \dots, G_d \in \R^{m \times m}$ such that, entrywise,
\begin{equation}\label{eq:ch2-G-conc}
G_\ell^{(N)} \approxP G_\ell, \qquad \ell = 1, \dots, d.
\end{equation}
With $G_{d+1}^{\alpha\beta} := 1$, their entries satisfy the backward recursion
\begin{equation}\label{eq:ch2-G-rec}
G_\ell = \Phi_\ell' \odot G_{\ell+1}, \qquad \ell = d, d-1, \dots, 1,
\end{equation}
where $\odot$ denotes the Hadamard (entrywise) product. Consequently,
\begin{equation}\label{eq:ch2-K0-conc}
K_0^{(N)} \approxP K := \sum_{\ell=0}^{d} G_{\ell+1} \odot \Phi_\ell
\end{equation}
entrywise (and hence $\st(K_0^{(N)}) = K$ $\Loeb$-almost surely).
\end{proposition}

\begin{proofsketch}
First consider the top hidden layer. The exact internal identity~\eqref{eq:ch2-gd-top} gives
\begin{equation}\label{eq:ch2-Gd-sum}
G_d^{(N),\alpha\beta} = \frac{1}{N} \sum_{i=1}^{N} W_{d,i}^2 \, \ns{\varphi}'(h_{d,i}^\alpha)\ns{\varphi}'(h_{d,i}^\beta).
\end{equation}
Since $W_d$ is internally independent of $\mathcal{F}_d$ and $\hE[W_{d,i}^2] = 1$, the internal conditional law-of-large-numbers calculation gives $G_d^{(N),\alpha\beta} \approxP \Phi_d^{\prime,\alpha\beta}$. Thus $G_d = \Phi_d'$.

For $1 \le \ell \le d-1$, take internal inner products in~\eqref{eq:ch2-gell-rec} to obtain
\begin{equation}\label{eq:ch2-Gell-quad}
G_\ell^{(N),\alpha\beta} = \frac{1}{N^2} (g_{\ell+1}^\alpha)^\top W_\ell \diag\bigl(\ns{\varphi}'(h_\ell^\alpha) \odot \ns{\varphi}'(h_\ell^\beta)\bigr) W_\ell^\top g_{\ell+1}^\beta.
\end{equation}
At the formal level of this proof sketch, transferred Gaussian conditioning (Lemma~\ref{lem:ch2-gauss-cond}, where $P_{\varphi_\ell}$ has standard rank at most $m$ inside $\hR^N$) suggests that the finite-rank projection terms are infinitesimal and the leading contribution is obtained by replacing $W_\ell$ in this quadratic form by a fresh internally independent Gaussian matrix $\widetilde{W}_\ell$. Conditional on the displayed $g$'s and $h$'s, the transferred elementary Gaussian quadratic-form identity gives
\begin{equation}\label{eq:ch2-Gell-fresh-exp}
\hE_{\widetilde{W}_\ell}\!\left[ \frac{1}{N^2} (g_{\ell+1}^\alpha)^\top \widetilde{W}_\ell \diag\bigl(\ns{\varphi}'(h_\ell^\alpha) \odot \ns{\varphi}'(h_\ell^\beta)\bigr) \widetilde{W}_\ell^\top g_{\ell+1}^\beta \right]
= G_{\ell+1}^{(N),\alpha\beta} \Phi_\ell^{\prime(N),\alpha\beta}.
\end{equation}
At unlimited width $N \in \hNinf$, the normalized internal quadratic form concentrates infinitesimally around this conditional expectation, so
\begin{equation}\label{eq:ch2-Gell-diff-inf}
G_\ell^{(N)} - \Phi_\ell^{\prime(N)} \odot G_{\ell+1}^{(N)} \approxP 0
\end{equation}
entrywise. Taking standard parts and proceeding by backward external induction through the standard finitely many layers $\ell = d, d-1, \dots, 1$ gives~\eqref{eq:ch2-G-conc} and~\eqref{eq:ch2-G-rec}. Finally, substituting the forward and backward shadows into the exact internal decomposition~\eqref{eq:ch2-layerwise-ntk}, with $t = 0$, gives~\eqref{eq:ch2-K0-conc}.
\end{proofsketch}

\begin{remark}[Rigor of the Proof Sketch]\label{rem:ch2-rigor-sketch}
The sketch suppresses the estimates needed to show that the finite-rank projection terms from Lemma~\ref{lem:ch2-gauss-cond} are infinitesimal in the fresh-Gaussian replacement and to justify the infinitesimal concentration of the normalized internal quadratic forms. At standard depth $d \in \N$, these steps can be made rigorous by applying transferred Gaussian conditioning recursively. The formal calculation above is the part that determines the standard shadow recursion.
\end{remark}

\begin{remark}[Other Architectures]\label{rem:ch2-other-arch}
The same strategy---identify the relevant internal ``state'' variables, use transferred Gaussian conditioning (Lemma~\ref{lem:ch2-gauss-cond}) to recover effective internal independence, then close a self-consistent standard-part recursion---extends beyond fully-connected networks (e.g.\ certain convolutional nets, residual networks in suitable regimes, and some attention-like blocks). The details depend strongly on the architecture.
\end{remark}

We have so far focused on optimization and training dynamics in the kernel regime. At unlimited width $N \in \hNinf$, when the standard limiting kernel $K$ is positive definite, the standard-part squared training loss converges exponentially to zero. This does not address the statistical properties of the trained predictor, particularly its generalization error on unseen data. We turn to this question in the next chapter.

\section{Exercises}\label{sec:ch2-exercises}

\begin{enumerate}
\item \textbf{Residual ODE.} Derive Proposition~\ref{prop:ch2-residual-ode} directly from the internal chain rule (by Transfer).

\item \textbf{General Loss Dissipation.} Prove Proposition~\ref{prop:ch2-risk-dissipation}. Then pass to standard parts in the fixed-kernel shadow ($K_t^{(N)} \approx K$) and show that, for any continuously differentiable standard pointwise loss $\ell$ (not necessarily squared) and any fixed standard symmetric positive semidefinite matrix $K$, the standard-part trajectories $f_\alpha(t)$ and $r_\alpha(t) = \partial_f \ell(f_\alpha(t), y^\alpha)$ satisfy
\begin{equation}\label{eq:ch2-ex-gen-loss}
\frac{d}{dt}\left[ \frac{1}{m} \sum_{\alpha=1}^{m} \ell\bigl(f_\alpha(t), y^\alpha\bigr) \right] = -\frac{1}{m^2} r(t)^\top K r(t) \le 0.
\end{equation}

\item \textbf{Fixed-Kernel Dynamics.} Verify Proposition~\ref{prop:ch2-fixed-kernel-exp} directly from the spectral theorem.

\item \textbf{One-Hidden-Layer NTK Formula.} At internal time $t \in \hR_{\ge 0}$ and unlimited width $N \in \hNinf$, write $a_t := a(t)$, $h_t(x) := W(t)x / \sqrt{n_0}$, and $f_t(x) = N^{-1/2} a_t^\top \ns{\varphi}(h_t(x))$. Show that the internal empirical NTK admits the exact decomposition
\begin{equation}\label{eq:ch2-ex-1layer-ntk}
K_t^{(N)}(x, x')
= \underbrace{\frac{1}{N} \ip{\ns{\varphi}(h_t(x))}{\ns{\varphi}(h_t(x'))}}_{\text{output-weight contribution}}
+ \underbrace{\frac{\ip{x}{x'}}{n_0} \frac{1}{N} \ip{a_t \odot \ns{\varphi}'(h_t(x))}{a_t \odot \ns{\varphi}'(h_t(x'))}}_{\text{feature-gradient term}}.
\end{equation}

\item \textbf{Deep NTK Recursion.} From the definition of $g_\ell^\alpha$ in~\eqref{eq:ch2-g-def}, derive the three boundary-aware internal identities in~\eqref{eq:ch2-layer-grads} and prove the exact decomposition~\eqref{eq:ch2-layerwise-ntk}. Verify the top-layer hyperfinite law-of-large-numbers calculation and the fresh-Gaussian internal expectation~\eqref{eq:ch2-Gell-fresh-exp} in the proof of Proposition~\ref{prop:ch2-backward-rec}. Use these calculations formally to obtain~\eqref{eq:ch2-Gell-diff-inf}, and then complete the backward external induction over standard $\ell = d, \dots, 1$ yielding $G_\ell^{(N)} \approxP G_\ell$ and $K_0^{(N)} \approxP K$.
\end{enumerate}
